\section{Lattice computations}
\label{app:lattices}

The seven integral lattices used in the manuscript are recorded here in
one place: the hyperbolic plane $U=\mathrm{II}_{1,1}$, the Lorentzian
root lattice $\La_{II}^{2,1}$ of the Borcherds--Kac--Moody algebra
$\autcor$, the K3 cohomology lattice $\La_{K3}=\mathrm{II}_{3,19}$, the
lattice-polarized K3 sublattice $\La_S$, the Mukai lattice
$\widetilde H(S,\Z)\simeq\mathrm{II}_{4,20}$, the formal dyonic charge
lattice $\Gamma_X^{\mathrm{form}}=\La_S\oplus\La_S$, and the
signature-$(3,2)$ lattice $\La^{3,2}$ that bridges $\Sp_4(\Z)$ to
$\SO_+(3,2)$ via $\PSp_4(\Z)\simeq\SO_+(\La^{3,2})$.  Each is given by
its Gram matrix or root data, signature, discriminant, and
automorphism group, with primary citations (Conway--Sloane
1999~\S2.5,~\S4,~\S16; Nikulin 1980 Theorem~1.1.2 and
Theorem~1.14.4; Mukai 1984; Gritsenko--Nikulin 1998
Part~II~\S2,~\S5).  Ranks, signatures, and the Cartan matrix of the
type-II simple roots are verified against the rational computation
\texttt{compute/verify\_lattice.py}.

This appendix is upstream of every chapter: load-bearing identities
about $\Delta_5$, $\autcor$, the orientation character $\nu_{\Delta_5}$,
and the Pfaffian wall normal form on
$W^{(2)}(\La_{II}^{2,1})$ refer back to the lattice data fixed
here, not to a sign or scaling chosen in the body.

\begin{convention}[Sign of the bilinear form; orientation of the
time-like cone]
\label{conv:lattice-signs}
The standard mathematical convention for the hyperbolic plane $U$ is
$U=(\Z^2,((0,1),(1,0)))$, even unimodular of signature $(1,1)$ and
discriminant $-1$ (Conway--Sloane 1999~\S2.5).  This appendix takes
this convention as primary.

The Borcherds--Gritsenko--Nikulin packaging used in
\S\ref{ssec:exceptional-iso} of the body and the exterior-square
model of $\Sp_4(\Z)$ employs the negated form
\[
\HypPlan=\Big(\Z^2,\begin{pmatrix}0&-1\\-1&0\end{pmatrix}\Big),
\]
so that the time-like cone $\Cone(\La^{2,1})_+\subset
\La^{2,1}\otimes\R$ has \emph{negative} square $(x,x)<0$.  The
relation is $\HypPlan=U(-1)$.  Every Gram matrix below is given in the
standard convention; a sign flip transports it to the body
convention without changing rank, discriminant absolute value, or
automorphism group.
\end{convention}

\subsection{The hyperbolic plane $U=\mathrm{II}_{1,1}$}
\label{app:lattice-U}

The hyperbolic plane is the unique even unimodular indefinite lattice
of signature $(1,1)$ (Conway--Sloane 1999~\S2.5; Nikulin 1980
Theorem~1.1.2).

\begin{definition}[$U=\mathrm{II}_{1,1}$]
\label{def:hyperbolic-plane}
\textnormal{[definitional, $\alpha$]}
Let $U:=\Z e\oplus\Z f$ with Gram matrix
\[
G_U=\begin{pmatrix}0&1\\1&0\end{pmatrix},
\qquad e^2=f^2=0,\quad e\cdot f=1.
\]
\end{definition}

\begin{verification}[Rank, signature, discriminant of $U$]
\label{ver:U-invariants}
\textnormal{[computed, $\gamma$]}
$\rk U=2$.  The eigenvalues of $G_U$ are $\pm1$, so the signature is
$(1,1)$.  The discriminant is
\[
\operatorname{disc} U=\det G_U=-1.
\]
$U$ is even (every $G_U$-diagonal entry is $0\equiv 0\pmod2$) and
unimodular ($|\det G_U|=1$).  These are the defining invariants of
$\mathrm{II}_{1,1}$ in the Conway--Sloane classification (Conway--Sloane
1999~\S2.5, Theorem~2.5.1).
\end{verification}

\begin{proposition}[$\mathrm{Aut}(U)$]
\label{prop:U-aut}
\textnormal{[proved, $\alpha$]}
The orthogonal group of $U$ is
\[
\mathrm{O}(U)=\{\pm\Id\}\times\langle\sigma\rangle
\simeq(\Z/2)^{2},
\qquad \sigma:e\leftrightarrow f.
\]
The proper subgroup $\mathrm{O}^+(U)\simeq\Z/2$ swaps $e$ and $f$ but
fixes the spinor norm.
\end{proposition}

\begin{proof}
An isometry $\phi\in\mathrm{O}(U)$ preserves the isotropic vectors of
$U$.  In rank $2$ over $\Z$, the isotropic primitive vectors are
exactly $\pm e,\pm f$, so $\phi$ permutes these four vectors.  The
permutations preserving the bilinear form give the four-element group
above (Conway--Sloane 1999~\S4.1).
\end{proof}

\begin{remark}[Standard sign convention versus body convention]
\label{rem:U-vs-HypPlan}
The body of the manuscript uses
$\HypPlan=U(-1)$
in the cusp polarization of \S\ref{ssec:exceptional-iso} and in the
exterior-square model of the automorphic group, where the time-like
cone of $\La^{2,1}$ has negative square.
The map
$U\to U(-1)$, $(e,f)\mapsto(e,f)$
is a bijection of underlying $\Z$-modules; only the bilinear form
flips sign.  Lattices that occur in the body
($\La^{2,1}$, $\La^{3,2}$) are presented in body coordinates so that
the Gram matrices match line-for-line with the manuscript text;
even-unimodular ambient lattices ($\La_{K3}$, $\widetilde H(S,\Z)$)
are presented in the standard Conway--Sloane convention.  Each
section flags its choice explicitly.
\end{remark}

\subsection{The Lorentzian root lattice $\La_{II}^{2,1}$}
\label{app:lattice-II21}

The Borcherds--Kac--Moody algebra $\autcor$ of
\cite[Theorem~2.1, Section~5.1.1, case $(t=1,II,\overline 0)$]{GNII}
is graded by an integral lattice of signature $(2,1)$.  This is the
lattice on which the type-II simple-root Cartan matrix is
$2(\Id_3-\J_3)$ where $\J_3$ is the all-ones matrix.

\begin{definition}[Ambient $\La^{2,1}=U(-1)\oplus\langle 2\rangle$]
\label{def:lambda-21-ambient}
\textnormal{[definitional, $\alpha$]}
Set
\[
\La^{2,1}:=U(-1)\oplus\langle 2\rangle
=\Z e\oplus\Z f\oplus\Z h,
\qquad e\cdot f=-1,\ e^2=f^2=0,\ h^2=2.
\]
The Gram matrix in the basis $(e,f,h)$ is
\[
G_{2,1}=\begin{pmatrix}0&-1&0\\ -1&0&0\\ 0&0&2\end{pmatrix},
\qquad \det G_{2,1}=-2,\qquad
\operatorname{disc}\La^{2,1}=-2.
\]
The body's lattice $\La^{2,1}=\HypPlan\oplus[2]$ in
$(f_2,f_3,f_{-2})$-coordinates uses exactly this Gram, with
$\HypPlan=U(-1)$ on the $(f_2,f_{-2})$-block and $[2]=\langle 2\rangle$
on the $f_3$-axis.
\end{definition}

\begin{definition}[$\La_{II}^{2,1}$ as the index-four type-II sublattice of $\La^{2,1}$]
\label{def:lambda-II-21}
\textnormal{[definitional, $\alpha$]}
The type-II sublattice
\[
\La_{II}^{2,1}\;:=\;
\{m f_2+l f_3+n f_{-2}\in\La^{2,1}\mid m\equiv 0\,(2),\ n\equiv 0\,(2)\}
\;\subset\;\La^{2,1}
\]
is the index-$4$ sublattice on which the type-II Cartan matrix
$2(\Id_3-\J_3)$ is the Gram of the simple roots
$(\delta_1,\delta_2,\delta_3)$ of
Proposition~\ref{prop:type-II-cartan}.  Abstractly it is the lattice
\[
\La_{II}^{2,1}\;\simeq\;U(4)\oplus\langle 2\rangle,\qquad
\operatorname{disc}\La_{II}^{2,1}=-32,
\]
matching the body terminology of Chapter~\ref{ch:pentadic} and the
type-II sublattice convention of Chapter~\ref{ch:view2-bkm}.
\end{definition}

\begin{verification}[Rank, signature, discriminant of $\La_{II}^{2,1}$]
\label{ver:II21-invariants}
\textnormal{[computed, $\gamma$]}
$\rk\La_{II}^{2,1}=3$.  The Gram of $\La_{II}^{2,1}$ in
$(\delta_1,\delta_2,\delta_3)$ of
Proposition~\ref{prop:type-II-cartan} is the Cartan matrix
$2(\Id_3-\J_3)$ of equation~\eqref{eq:bkm-cartan-matrix}; its
determinant is
\[
\det\bigl(2(\Id_3-\J_3)\bigr)=-32.
\]
The hyperbolic-plus-rank-one block of the ambient $\La^{2,1}$ has
signature $(2,1)$; the index-$4$ sublattice $\La_{II}^{2,1}$ inherits
this signature, so
\[
\sigma(\La_{II}^{2,1})=(2,1).
\]
The discriminant group $(\La_{II}^{2,1})^*/\La_{II}^{2,1}$ has order
$32$ as required.  The lattice $\La_{II}^{2,1}$ is even because every
$\delta_i^2=2$ and the off-diagonal couplings $(\delta_i,\delta_j)=-2$
are integral.  Switching to the standard hyperbolic-plane convention
$U=((0,1),(1,0))$ on the ambient $\La^{2,1}$ presents
$\La^{2,1}\simeq U\oplus\langle -2\rangle$ and
$\La_{II}^{2,1}\simeq U(4)\oplus\langle 2\rangle$; rank, signature,
discriminant group order, automorphism group, and Cartan datum are
unchanged.
\end{verification}

\begin{remark}[Identification with the body lattice]
\label{rem:II21-vs-La21}
The body uses
$\La^{2,1}=\HypPlan\oplus[2]
=\Z f_2\oplus\Z f_3\oplus\Z f_{-2}$
with Gram in $(f_2,f_3,f_{-2})$
\[
\begin{pmatrix}0&0&-1\\ 0&2&0\\ -1&0&0\end{pmatrix},
\]
matching $G_{2,1}$ of Definition~\ref{def:lambda-21-ambient} up to
the basis swap $e\mapsto f_2$, $f\mapsto f_{-2}$, $h\mapsto f_3$.
The type-II sublattice $\La_{II}^{2,1}\subset\La^{2,1}$ of
Definition~\ref{def:lambda-II-21} is the index-$4$ sublattice on
which the simple roots $\delta_1=2f_2-f_3$, $\delta_2=2f_{-2}-f_3$,
$\delta_3=f_3$ of Proposition~\ref{prop:type-II-cartan} live; it is
the even Lorentzian root lattice of $\autcor$.
\end{remark}

\begin{proposition}[Type-II simple-root Cartan]
\label{prop:type-II-cartan}
\textnormal{[proved, $\gamma$]}
In the body coordinates fixed at the start of the Weyl-chamber
section, set
\[
\delta_1=2f_2-f_3,\qquad
\delta_2=2f_{-2}-f_3,\qquad
\delta_3=f_3.
\]
Then $\delta_i^2=2$ for $i=1,2,3$ and
\[
\bigl((\delta_i,\delta_j)\bigr)_{i,j=1,2,3}
=\begin{pmatrix} 2&-2&-2\\ -2&2&-2\\ -2&-2&2\end{pmatrix}
=2(\Id_3-\J_3),
\]
where $\J_3$ is the $3\times 3$ all-ones matrix.  Each $\delta_i$ is
type-II in the sense that $(\delta_i,\La^{2,1})=2\Z$.
\end{proposition}

\begin{proof}
Direct computation in the body Gram.  The companion script
\texttt{compute/verify\_lattice.py} at lines 92--97 verifies the matrix
$2(\Id_3-\J_3)$ from the same data; the print output records
\texttt{Type-II simple-root Gram matrix} equal to the displayed Cartan.
The type-II divisibility statement follows because each $\delta_i$
pairs against $\La^{2,1}$ with even values: $(\delta_i,f_2)$,
$(\delta_i,f_3)$, $(\delta_i,f_{-2})\in 2\Z$ for each $i$.
This is the GN root datum of \cite[Section 5.1.1, case
$(1,II,\overline0)$]{GNII}.
\end{proof}

\begin{proposition}[Weyl vector $\rho$]
\label{prop:weyl-vector-II21}
\textnormal{[proved, $\gamma$]}
The Weyl vector
\[
\rho:=\tfrac12(\delta_1+\delta_2+\delta_3)
=f_2-\tfrac12 f_3+f_{-2}\in\La^{2,1}\otimes\Q
\]
satisfies $(\rho,\delta_i)=-1$ for $i=1,2,3$ and
\[
(\rho,\rho)=-\tfrac32,\qquad \rho\in\Cone(\La^{2,1})_+.
\]
In particular $\rho$ lies in the time-like cone, so the chamber
$\Poly_{II}=\{x\mid(x,\delta_i)\le 0,\,i=1,2,3\}$ has finite hyperbolic
volume.
\end{proposition}

\begin{proof}
The simple-root condition $(\rho,\delta_i)=-(\delta_i^2)/2=-1$
follows from Kac~\cite[\S2.5, Lemma~2.5]{KAC:1} for any choice of
simple roots of square $2$; for the GN type-II datum see
\cite[(2.7)]{GNII}.  The script
\texttt{compute/verify\_lattice.py} asserts $\rho=(1,-1/2,1)$ in
$(f_2,f_3,f_{-2})$ coordinates (line~99) and
$(\rho,\delta_i)=-1$ at line~100.  Direct computation in the body
Gram of Definition~\ref{def:lambda-II-21}:
\[
(\rho,\rho)
=2(f_2,f_{-2})+\tfrac14(f_3,f_3)
=2(-1)+\tfrac14(2)=-\tfrac32.
\]
Since $(\rho,\rho)<0$, $\rho\in\Cone(\La^{2,1})_+$.
\end{proof}

\begin{proposition}[Reflection group; Nikulin's theorem]
\label{prop:nikulin-II21-reflection}
\textnormal{[proved, $\alpha$; Nikulin 1980 Theorem~1.4.1]}
Let $W^{(2)}(\La^{2,1})$ be the subgroup of $\Orth(\La^{2,1})$
generated by reflections in primitive roots of square $2$.  Then
\[
\Orth(\La^{2,1})_+=W^{(2)}(\La^{2,1}),
\]
where the subscript $+$ denotes the index-two subgroup preserving the
time-like cone $\Cone(\La^{2,1})_+$.  Equivalently, every isometry of
$\La^{2,1}$ preserving the chosen positive cone is a product of
square-$2$ reflections.
\end{proposition}

\begin{proof}
Nikulin~\cite[Theorem~1.4.1]{NIKULIN} establishes for the rank-three
hyperbolic lattice $\HypPlan\oplus[2]$ the equality
$\Orth(\La^{2,1})_+=W^{(2)}(\La^{2,1})$; see also
Gritsenko--Nikulin~\cite[Proposition 1.5]{GN}.  The body of the
manuscript invokes this at line~5408--5410.
\end{proof}

\begin{verification}[Type-II reflection orbit]
\label{ver:type-II-reflection-orbit}
\textnormal{[computed, $\gamma$]}
For each $\delta\in\{\delta_1,\delta_2,\delta_3\}$, the reflection
$s_\delta:x\mapsto x-(x,\delta)\delta$ satisfies
$s_\delta(\delta)=-\delta$ and preserves the Cartan matrix on the
$\delta$-orbit (\texttt{compute/verify\_lattice.py} lines~103--108).
Concretely, $s_{\delta_1}(\delta_1)=-\delta_1$,
$s_{\delta_1}(\delta_2)=\delta_2+2\delta_1$, etc.; the bilinear form
on $\{\delta_1,\delta_2,\delta_3\}$ is reflection-invariant by direct
matrix computation.
\end{verification}

\begin{remark}[The two type-$2$ divisibility classes]
\label{rem:type-I-type-II}
Each $\delta\in\La^{2,1}$ with $\delta^2=2$ has divisibility
$(\delta,\La^{2,1})\in\{\Z,2\Z\}$.  The two cases are recorded as
\emph{type I} ($(\delta,\La^{2,1})=\Z$) and \emph{type II}
($(\delta,\La^{2,1})=2\Z$); both classes generate index-finite normal
subgroups of $W^{(2)}(\La^{2,1})$:
\[
[W^{(2)}(\La^{2,1}):W^{(2)}(\La^{2,1}_I)]=2,\qquad
[W^{(2)}(\La^{2,1}):W^{(2)}(\La^{2,1}_{II})]=6,
\]
following \cite[Sections 4--5]{GN}.  The denominator algebra
$\autcor$ uses the type-II reflection chamber.
\end{remark}

\subsection{The K3 lattice $\La_{K3}=\mathrm{II}_{3,19}$}
\label{app:lattice-K3}

For $S$ a complex K3 surface, the integral cohomology
$H^2(S,\Z)$ with cup-product pairing is the unique even unimodular
indefinite lattice of signature $(3,19)$
(Conway--Sloane 1999~\S16.4; Nikulin 1980 Theorem~1.1.4).

\begin{definition}[$\La_{K3}=\mathrm{II}_{3,19}$]
\label{def:K3-lattice}
\textnormal{[definitional, $\alpha$]}
The K3 cohomology lattice is
\[
\La_{K3}:=H^2(S,\Z)\simeq -E_8\oplus -E_8\oplus U\oplus U\oplus U
\equiv -2E_8\oplus 3U,
\]
where $E_8$ is the rank-$8$ root lattice with the standard positive-
definite Cartan pairing (Conway--Sloane 1999~\S4.8) and $-E_8$
denotes $E_8$ with bilinear form negated.
\end{definition}

\begin{verification}[Rank, signature, discriminant of $\La_{K3}$]
\label{ver:K3-invariants}
\textnormal{[computed, $\alpha$]}
\[
\rk\La_{K3}=8+8+2+2+2=22,
\qquad
\sigma(\La_{K3})=(3,19).
\]
The three $U$-summands contribute signature $(3,3)$ and the two
$-E_8$-summands contribute signature $(0,16)$; total $(3,19)$.  The
discriminant is
\[
\operatorname{disc}\La_{K3}=
(\det -E_8)^2\cdot(\det U)^3=1\cdot(-1)^3=-1,
\]
so $|\operatorname{disc}|=1$ and the lattice is unimodular
(Conway--Sloane 1999~\S16.4, Theorem~16.4.1).  Even unimodularity
follows because each summand is even unimodular.
\end{verification}

\begin{proposition}[Aut and Hodge data on $\La_{K3}$]
\label{prop:K3-aut}
\textnormal{[proved, $\alpha$]}
The orthogonal group $\Orth(\La_{K3})$ acts transitively on the set
of primitive vectors of any fixed square (Eichler 1952;
Nikulin~\cite[Theorem 1.14.4]{Nikulin1980}).  For a polarized K3 with
period $\omega\in H^{2,0}(S)\subset\La_{K3}\otimes\C$, the
transcendental and algebraic sublattices fit into the orthogonal
splitting
\[
\La_{K3}=\NS(S)\oplus T(S),
\]
where $\NS(S)\subset H^{1,1}(S,\Z)$ is the N\'eron--Severi
(algebraic) sublattice of signature $(1,\rho-1)$ and $T(S)$ is the
transcendental lattice of signature $(2,20-\rho)$, with
$\rho=\operatorname{rk}\NS(S)$ the Picard rank
(Beauville--Bourguignon--Demazure 1985~\S6;
\cite[\S2]{Nikulin1980}).
\end{proposition}

\begin{proof}
The transitivity statement is Eichler's theorem on indefinite
unimodular lattices, sharpened by
Nikulin~\cite[Theorem~1.14.4]{Nikulin1980}: for an even unimodular
indefinite lattice $L$ of rank $\ge 3$, $\Orth(L)$ acts transitively
on primitive vectors of fixed square in each squared-class.
Decomposition by Hodge type partitions $\La_{K3}\otimes\C$ into
$H^{2,0}\oplus H^{1,1}\oplus H^{0,2}$; intersecting with $\La_{K3}$
gives the transcendental--algebraic splitting at the integral level.
\end{proof}

\subsection{The lattice-polarized K3 sublattice $\La_S$}
\label{app:lattice-Lambda-S}

A lattice polarization is a primitive embedding of an integral
lattice $\La_S$ into $\La_{K3}$ whose orthogonal complement carries
period structure.

\begin{definition}[Lattice polarization $\La_S\hookrightarrow\La_{K3}$]
\label{def:lattice-polarization}
\textnormal{[definitional, $\alpha$]}
A \emph{lattice polarization} of signature $(1,\rho-1)$ on a K3
surface $S$ is a primitive embedding
\[
\La_S\hookrightarrow\La_{K3}=\mathrm{II}_{3,19},\qquad
\sigma(\La_S)=(1,\rho-1),\quad 1\le\rho\le 19,
\]
where $\La_S\subset H^{1,1}(S,\Z)$ contains an ample class.  The
orthogonal complement
\[
\La_S^\perp\subset\La_{K3},\qquad \sigma(\La_S^\perp)=(2,20-\rho),
\]
is the lattice on which the period map of the family
$\mathcal F_{\La_S}\to\mathcal D_{\La_S^\perp}$ is defined
(Dolgachev~1996; Nikulin 1980~Proposition~1.6.1).
\end{definition}

\begin{verification}[Numerics of $\La_S$ and $\La_S^\perp$]
\label{ver:lambda-S-invariants}
\textnormal{[computed, $\alpha$]}
\[
\rk\La_S=\rho,\quad \sigma(\La_S)=(1,\rho-1),\quad
\rk\La_S^\perp=22-\rho,\quad \sigma(\La_S^\perp)=(2,20-\rho).
\]
The discriminant pairs:
\[
\operatorname{disc}\La_S\cdot\operatorname{disc}\La_S^\perp
=\operatorname{disc}\La_{K3}=-1,
\]
so $\operatorname{disc}\La_S^\perp=-\operatorname{disc}\La_S^{-1}\cdot
\det A^2$ for $A$ the change-of-basis matrix realizing the primitive
embedding (Nikulin 1980~Corollary~1.6.2).  In particular
$\La_S^\perp/\La_S^*$ and $\La_{K3}/\La_S\oplus\La_S^\perp$ are
isomorphic as discriminant groups.
\end{verification}

\begin{proposition}[Nikulin embedding criterion]
\label{prop:nikulin-embedding}
\textnormal{[proved; Nikulin 1980 Theorem~1.14.4]}
A primitive embedding $\La_S\hookrightarrow\La_{K3}$ exists if and
only if there is an isomorphism
\[
(A_{\La_S},q_{\La_S})\simeq(A_{\La_S^\perp},-q_{\La_S^\perp})
\]
of finite quadratic forms, where $A_L:=L^*/L$ is the discriminant
group and $q_L$ its quadratic form (Nikulin 1980 Definition~1.3.1
and Theorem~1.14.4; Conway--Sloane 1999~\S15.5).  The embedding is
unique up to $\Orth(\La_{K3})$ when
\[
\rk\La_S+\ell(A_{\La_S})\le 22-\rk\La_S
\]
and the discriminant form satisfies the Nikulin uniqueness criterion
(Nikulin 1980 Theorem~1.14.4).
\end{proposition}

\begin{proof}
This is the gluing-via-discriminant-form theorem
(\cite[\S1.14]{Nikulin1980}; \cite[\S15.5]{ConwaySloane}).  An
embedding into $\La_{K3}$ exists iff the discriminant form on
$\La_S^*/\La_S$ is realized by the orthogonal complement; uniqueness
is the Nikulin orbit-uniqueness statement under the corank
inequality.
\end{proof}

\begin{remark}[Mirror lattice convention]
\label{rem:mirror-lattice-convention}
The Aspinwall--Morrison mirror to a $\La_S$-polarized K3 family is the
$\check\La_S$-polarized family with
\[
\check\La_S=U\oplus(\La_S^\perp\ominus U),
\]
the splitting depending on a choice of isotropic primitive
$U\hookrightarrow\La_S^\perp$ (Dolgachev 1996~\S5; Aspinwall 1997).
The two factors $\La_S$ and $\La_S^\perp\ominus U$ play the roles of
algebraic--symplectic exchange under mirror symmetry.  No mirror
identification is asserted at this lattice level; the assertion that
$\Delta_5^{-2}$ is the discriminant of the mirror period family
(View~\textcircled{4}) is conjectural and not used in
Chapter~\ref{chap:dirac-igusa-pro-object}.
\end{remark}

\subsection{The Mukai lattice $\widetilde H(S,\Z)=\mathrm{II}_{4,20}$}
\label{app:lattice-Mukai}

For $S$ a complex K3 surface, the full integral cohomology
$H^*(S,\Z)$ with the Mukai pairing is an even unimodular lattice of
signature $(4,20)$ (Mukai 1984~Theorem~1.5; Mukai 1987).

\begin{definition}[Mukai lattice $\widetilde H(S,\Z)$]
\label{def:mukai-lattice-appA}
\textnormal{[definitional, $\alpha$]}
The Mukai lattice of $S$ is
\[
\widetilde H(S,\Z)
:=H^0(S,\Z)\oplus H^2(S,\Z)\oplus H^4(S,\Z),
\qquad
\widetilde H(S,\Z)\simeq U\oplus U\oplus U\oplus U\oplus -2E_8
\equiv 4U\oplus 2(-E_8),
\]
graded by even cohomological degree.  The Mukai pairing is
\[
\langle (r,D,s),(r',D',s')\rangle_{\mathrm{Muk}}
:=D\cdot D'-rs'-r's
\]
for $(r,D,s),(r',D',s')\in
\Z\oplus H^2(S,\Z)\oplus\Z$
(Mukai 1984~Theorem~1.5).  The pairing of an $H^0$-class with an
$H^4$-class carries a minus sign relative to the cup product, so that
the rank-$4$ component of $H^*(S,\Z)\otimes H^*(S,\Z)$ contributes
$-rs'-r's$.
\end{definition}

\begin{verification}[Rank, signature, discriminant of
$\widetilde H(S,\Z)$]
\label{ver:mukai-invariants}
\textnormal{[computed, $\alpha$]}
\[
\rk\widetilde H(S,\Z)
=1+22+1=24,
\qquad
\sigma(\widetilde H(S,\Z))=(4,20).
\]
The four $U$-summands give signature $(4,4)$ and the $-2E_8$ pieces
give signature $(0,16)$, total $(4,20)$.  Discriminant
\[
\operatorname{disc}\widetilde H(S,\Z)
=(\det U)^4\cdot(\det -E_8)^2
=(-1)^4\cdot 1=1,
\]
so $\widetilde H(S,\Z)$ is even unimodular, denoted
$\mathrm{II}_{4,20}$ (Conway--Sloane 1999~\S16.4).  The body uses the
notation $\La_S=\widetilde H(S,\Z)$ at line~4589 and~4714 of
\texttt{main.tex}; this is consistent with the convention here.
\end{verification}

\begin{proposition}[Mukai vector]
\label{prop:mukai-vector}
\textnormal{[proved, $\alpha$; Mukai 1987]}
For $F$ a coherent sheaf on $S$, the Mukai vector is
\[
v(F):=\ch(F)\sqrt{\td(S)}
=(\rk F,\,c_1(F),\,\ch_2(F)+\rk F\cdot[\mathrm{pt}])
\in\widetilde H(S,\Z),
\]
where $\sqrt{\td(S)}=(1,0,1)$ in Mukai notation on a K3 surface
(\cite[Theorem~3]{Mukai1987}).  The Mukai pairing on
$\widetilde H(S,\Z)$ recovers the Euler characteristic of $\Hom$:
\[
\langle v(F),v(G)\rangle_{\mathrm{Muk}}
=-\chi(F,G)
=-\sum_i(-1)^i\dim\Ext^i(F,G).
\]
\end{proposition}

\begin{proof}
Mukai~\cite[Theorem~3]{Mukai1987} computes $\sqrt{\td(S)}$ for K3
surfaces.  The Hirzebruch--Riemann--Roch formula gives
$\chi(F,G)=\int_S\ch(F^\vee)\ch(G)\td(S)
=-\langle v(F),v(G)\rangle_{\mathrm{Muk}}$ once the sign of the
$H^0$--$H^4$ pairing is taken into account; see Mukai 1984~\S2.
\end{proof}

\begin{remark}[Sign of the Mukai pairing on rank--$0$ components]
\label{rem:mukai-rank-0-sign}
The Mukai pairing
$\langle(r,D,s),(r',D',s')\rangle_{\mathrm{Muk}}
=D\cdot D'-rs'-r's$
restricts to the cup-product pairing on the middle $H^2$-component
\[
\langle(0,D,0),(0,D',0)\rangle_{\mathrm{Muk}}=D\cdot D',
\]
and to the negative cup-product on the rank--Euler component
\[
\langle(r,0,s),(r',0,s')\rangle_{\mathrm{Muk}}=-rs'-r's.
\]
The relative minus sign is forced by Hirzebruch--Riemann--Roch
\cite[Theorem~5.1.1]{HuybrechtsLehn}: $-\chi(F,G)$ contains a
$-\rk F\cdot\chi(\O,G)-\rk G\cdot\chi(\O,F)$ term coming from
$\sqrt{\td(S)}=(1,0,1)$.  An apparent attack route is to negate the
$H^0$--$H^4$ pairing or to drop the minus sign; this would convert
the Mukai lattice to an isotropic-rank-$2$ unimodular form not of
signature $(4,20)$, in disagreement with the Mukai 1984 classification
(\cite[Theorem~1.5]{Mukai1984}) and the unique even unimodular
type II$_{4,20}$.  The sign as stated is forced by these
requirements.
\end{remark}

\begin{proposition}[Aut of the Mukai lattice]
\label{prop:mukai-aut}
\textnormal{[proved, $\alpha$]}
$\Orth(\widetilde H(S,\Z))$ acts transitively on the set of primitive
vectors of any fixed square (Eichler 1952; Nikulin 1980
Theorem~1.14.4).  In particular, the Mukai vector $v=(1,0,1)$
(corresponding to the structure sheaf $\mathcal O_S$) and the
generator $(0,0,1)$ of $H^4(S,\Z)\cap\widetilde H(S,\Z)$ are
$\Orth(\widetilde H(S,\Z))$-conjugate to any other primitive vector
of the same square.  The derived autoequivalence group of $D^b(S)$
maps to $\Orth(\widetilde H(S,\Z))$ via the action on Mukai vectors
(Mukai 1987~\S2; Orlov 1997).
\end{proposition}

\begin{proof}
Eichler's theorem (\cite[Theorem~1.14.4]{Nikulin1980}) gives the
transitivity statement.  Mukai~\cite[\S2]{Mukai1987} constructs the
representation of derived autoequivalences on
$\widetilde H(S,\Z)$ via Fourier--Mukai kernels.
\end{proof}

\subsection{The formal dyonic charge lattice
$\Gamma_X^{\mathrm{form}}=\La_S\oplus\La_S$}
\label{app:lattice-formal-dyonic}

The formal additive dyonic lattice on $X=S\times E$ packages the
electric--magnetic charge data on which the Gram-shadow projection
$\Pi_X$ is defined.  Throughout $\La_S=\widetilde H(S,\Z)$ in the
sense of \S\ref{app:lattice-Mukai}.

\begin{definition}[$\Gamma_X^{\mathrm{form}}$ and its Gram pairing]
\label{def:gamma-X-form}
\textnormal{[definitional, $\alpha$; from
Definition~\ref{def:mukai-lattice}]}
For $X=S\times E$ with $S$ a K3 surface and $E$ an elliptic curve,
\[
\Gamma_X^{\mathrm{form}}:=\La_S\oplus\La_S,\qquad
c=(Q,P)\in\Gamma_X^{\mathrm{form}},\quad Q,P\in\La_S.
\]
The Gram pairing
$\langle\,,\,\rangle_{\Gamma}:\Gamma_X^{\mathrm{form}}
\times\Gamma_X^{\mathrm{form}}\to\Z$ is inherited from the Mukai
pairing on each factor:
\[
\langle (Q,P),(Q',P')\rangle_\Gamma
:=\langle Q,Q'\rangle_{\mathrm{Muk}}
+\langle P,P'\rangle_{\mathrm{Muk}}.
\]
The Gram-index projection
\[
\Pi_X:\Gamma_X^{\mathrm{form}}\to\Gamma_{\mathrm{gram}}=\Z^3,\qquad
(Q,P)\mapsto\Big(\tfrac{Q^2}{2},\,Q\cdot P,\,\tfrac{P^2}{2}\Big),
\]
records the genus-two Gram triple of the dyonic class.  This map is
quadratic, not additive: the obstruction is the symmetric bilinear
$2$-cocycle $B$ of Lemma~\ref{lem:gram-additivity-defect}.
\end{definition}

\begin{verification}[Rank, signature of $\Gamma_X^{\mathrm{form}}$]
\label{ver:gamma-X-form-invariants}
\textnormal{[computed, $\alpha$]}
\[
\rk\Gamma_X^{\mathrm{form}}=2\cdot\rk\La_S=48,
\qquad
\sigma(\Gamma_X^{\mathrm{form}})=2\cdot(4,20)=(8,40).
\]
The discriminant is
\[
\operatorname{disc}\Gamma_X^{\mathrm{form}}
=(\operatorname{disc}\La_S)^2=1,
\]
so $\Gamma_X^{\mathrm{form}}$ is even unimodular.  In Conway--Sloane
notation it is $\mathrm{II}_{8,40}$ (Conway--Sloane 1999~\S16.4).
\end{verification}

\begin{remark}[Relation to the K3$\times$E Mukai geometry]
\label{rem:gamma-X-vs-K3E}
Under the K\"unneth decomposition for $X=S\times E$, the integral
even cohomology decomposes as
\[
H^{\mathrm{ev}}(X,\Z)
=H^*(S,\Z)\otimes H^*(E,\Z)
=H^*(S,\Z)\otimes(\Z[1]\oplus\Z[\omega_E]),
\]
where $\omega_E$ is the fundamental class of $E$.  Every Mukai vector
on $X$ in the formal even sector decomposes as
$v_X=P\otimes 1_E+Q\otimes\omega_E$ with $Q,P\in\La_S$, and
$\Gamma_X^{\mathrm{form}}=\La_S\oplus\La_S$ is the additive lattice
on which the $(P,Q)$-grading lives
(Definition~\ref{def:gram-index-lattice}).
The Mukai pairing on $\widetilde H(X,\Z)$ pairs against the elliptic
$U$-summand $H^0(E,\Z)\oplus H^2(E,\Z)\simeq U$, giving
\[
\langle v_X,v_X'\rangle_{\widetilde H(X)}
=Q\cdot P'+Q'\cdot P,
\]
where each summand is the Mukai pairing on $\La_S$.  The resulting
genus-two Gram triple $(Q^2/2,Q\cdot P,P^2/2)$ is the Gram-index
data $\Pi_X(Q,P)$ used in the Borcherds product
(Theorem~\ref{thm:d6-d2-d0-dictionary-appendixC}).
\end{remark}

\begin{lemma}[Two orthogonal hyperbolic planes inside $\La_S$]
\label{lem:two-hyperbolic-planes-appA}
\textnormal{[proved, $\alpha$]}
$\La_S=\widetilde H(S,\Z)\simeq 4U\oplus 2(-E_8)$ contains four
mutually orthogonal copies of $U$.  In particular, two of them give
\[
U_1\oplus U_2\subset\La_S,\qquad U_i\simeq U,\ U_1\perp U_2,
\]
which is the hypothesis of the formal primitive lift
Lemma~\ref{lem:primitive-lift-every-gram-triple}: every
Gram triple $(n,l,m)\in\Z^3$ has a primitive saturated formal lift
$c=(Q,P)\in\Gamma_X^{\mathrm{form}}$ with $\Pi_X(Q,P)=(n,l,m)$.
\end{lemma}

\begin{proof}
By Definition~\ref{def:mukai-lattice}, $\La_S\simeq 4U\oplus 2(-E_8)$
contains four mutually orthogonal $U$-summands.  Selecting any two of
them gives the displayed sublattice, which is itself isomorphic to
$U\oplus U\simeq\mathrm{II}_{2,2}$ as a sublattice of $\La_S$.  The
Gram-triple lift then follows from
Lemma~\ref{lem:primitive-lift-every-gram-triple}: take
$Q=e_1+nf_1$, $P=lf_1+e_2+mf_2$ in $U_1\oplus U_2$ to obtain
$\Pi_X(Q,P)=(n,l,m)$.
\end{proof}

\subsection{The signature-$(3,2)$ lattice $\La^{3,2}$ and the
$\Sp_4\simeq\Spin(3,2)$ bridge}
\label{app:lattice-32}

The exceptional isomorphism $\PSp_4(\R)\simeq\SO_+(3,2)$ identifies
the Siegel modular variety with the fourth-orthogonal modular variety
of an integral lattice $\La^{3,2}$ of signature $(3,2)$
(\cite[Section~3]{GN}; Borcherds 1995~\S5).  The Igusa cusp form
$\Delta_5$ lives on the Siegel side; the Borcherds product expansion
lives on the orthogonal side.

\begin{definition}[$\La^{3,2}$ in the exterior-square model]
\label{def:lambda-32}
\textnormal{[definitional, $\alpha$]}
Let $\La^4=\Z e_1\oplus\Z e_2\oplus\Z e_3\oplus\Z e_4$ and consider
the exterior square $\wedge^2\La^4$ with the symmetric pairing
\[
u\wedge v\cdot u'\wedge v'=
[\det(u,v,u',v')]_{e_1\wedge e_2\wedge e_3\wedge e_4}\in\Z.
\]
Then $(\wedge^2\La^4,(\,,\,))\simeq\mathrm{II}_{3,3}$ as an even
unimodular lattice of signature $(3,3)$.  Fix
$q_0:=e_1\wedge e_3+e_2\wedge e_4\in\wedge^2\La^4$ and define
\[
\La^{3,2}:=q_0^\perp\subset\wedge^2\La^4.
\]
\end{definition}

\begin{verification}[Rank, signature, discriminant of $\La^{3,2}$]
\label{ver:lambda-32-invariants}
\textnormal{[computed, $\alpha$; verified by
\texttt{compute/verify\_lattice.py}]}
$\rk\wedge^2\La^4=\binom{4}{2}=6$.  Removing one isotropic ray from
$q_0$ gives $\rk\La^{3,2}=5$.  The \texttt{expected\_gram} array of
\texttt{compute/verify\_lattice.py} at lines~82--88 is the explicit Gram
matrix of $\La^{3,2}$ in the basis $(f_1,f_2,f_3,f_{-2},f_{-1})$:
\[
G_{3,2}=\begin{pmatrix}
0&0&0&0&-1\\
0&0&0&-1&0\\
0&0&2&0&0\\
0&-1&0&0&0\\
-1&0&0&0&0
\end{pmatrix},
\]
where $f_1=e_1\wedge e_2$, $f_2=e_2\wedge e_3$,
$f_3=e_1\wedge e_3-e_2\wedge e_4$, $f_{-2}=e_4\wedge e_1$,
$f_{-1}=e_4\wedge e_3$.  Direct computation: the matrix is
antidiagonal except for the central $h^2=2$ entry, with off-diagonal
entries $-1,-1,-1,-1$.  Expansion gives
\[
\det G_{3,2}=2\cdot(-1)^4\cdot\operatorname{sgn}(1,2,3,4\mapsto4,3,2,1)
=2\cdot 1\cdot(+1)=+2,
\]
since the permutation $(1,2,3,4)\mapsto(4,3,2,1)$ has six inversions
and even parity.  Therefore
\[
\sigma(\La^{3,2})=(3,2),\qquad
\operatorname{disc}\La^{3,2}=2.
\]
The script asserts
\texttt{gram == expected\_gram} at line~89.
\end{verification}

\begin{proposition}[$\La^{3,2}\simeq U\oplus U\oplus\langle 2\rangle$
in the manuscript convention; $\La^{3,2}\simeq U(-1)\oplus
U(-1)\oplus\langle 2\rangle$ in the standard convention]
\label{prop:lambda-32-decomposition}
\textnormal{[proved, $\alpha$]}
In the body convention of the exterior-square model,
\[
\La^{3,2}\simeq\HypPlan\oplus\HypPlan\oplus[2],
\]
with $\HypPlan=U(-1)$ and $[2]$ the rank-$1$ lattice with Gram
$(2)$.  The Gram presentation displayed in
Verification~\ref{ver:lambda-32-invariants} arises from the
$f$-basis: $(f_1,f_{-1})$ pairs as a copy of $\HypPlan$
($(f_1,f_{-1})=-1$), $(f_2,f_{-2})$ pairs as a second copy of
$\HypPlan$ ($(f_2,f_{-2})=-1$), and $f_3$ contributes
$(f_3,f_3)=2$.
\end{proposition}

\begin{proof}
The pairings $(f_1,f_{-1})=-1$, $(f_2,f_{-2})=-1$, $(f_3,f_3)=2$ are
read off from $G_{3,2}$.  The orthogonal decomposition is direct:
$\Z f_1+\Z f_{-1}\perp\Z f_2+\Z f_{-2}\perp\Z f_3$ inside
$\La^{3,2}$ because every cross-pair vanishes in $G_{3,2}$.  The
identification with $\HypPlan\oplus\HypPlan\oplus[2]$ follows
because each rank-$2$ block has Gram $((0,-1),(-1,0))=\HypPlan$
and the rank-$1$ block has Gram $(2)=[2]$.
\end{proof}

\begin{proposition}[$\PSp_4(\Z)\simeq\SO_+(\La^{3,2})$]
\label{prop:psp4-iso-so32}
\textnormal{[proved, $\alpha$; exterior-square model of $\Sp_4(\Z)$]}
The exterior-square representation
$\wedge^2:\Sp_4(\Z)/\{\pm\Id_4\}\to\SO_+(\La^{3,2})\simeq
\Orth(\La^{3,2})_+/\{\pm\Id_5\}$
is an isomorphism.  Concretely, the integral generators
\[
g_0=\begin{pmatrix}0&\Id_2\\-\Id_2&0\end{pmatrix},\quad
g_M=\begin{pmatrix}\Id_2&M\\ 0&\Id_2\end{pmatrix},\quad
g_A=\begin{pmatrix}{}^tA^{-1}&0\\ 0&A\end{pmatrix}
\quad(A\in\GL_2(\Z))
\]
of $\Sp_4(\Z)$ map under $\wedge^2$ to integral isometries of
$\La^{3,2}$, displayed as the three explicit $5\times 5$ matrices
$\wedge^2(g_0)$, $\wedge^2(g_M)$, $\wedge^2(g_A)$ in the body's
exterior-square model.  The full image is $\SO_+(\La^{3,2})$.
The Siegel space $\Hpl_2$ embeds into the type-IV upper half-space
$\Hpl^{\IV}_+\subset\Proj(\La^{3,2}\otimes\C)$ via
\[
Z=\begin{pmatrix}z_1&z_2\\ z_2&z_3\end{pmatrix}
\mapsto
{}^t(z_2^2-z_1z_3,\,z_3,\,z_2,\,z_1,\,1)z_0,
\]
and the action of $\Sp_4(\Z)$ on $\Hpl_2$ corresponds to the action
of $\SO_+(\La^{3,2})$ on $\Hpl^{\IV}_+$.
\end{proposition}

\begin{proof}
The Pfaffian pairing on $\wedge^2\La^4$ is preserved by
$\wedge^2\SL_4$.  The subgroup fixing $q_0$ is $\Sp_4(\Z)$ by
definition of the symplectic form
$B_{q_0}(x,y)e_1\wedge e_2\wedge e_3\wedge e_4=-x\wedge y\wedge q_0$.
The kernel on $q_0^\perp=\La^{3,2}$ is $\{\pm\Id_4\}$.  The
exceptional Lie-algebra isomorphism $\mathfrak{sp}_4(\R)\simeq\mathfrak{so}(3,2)$
extends to the integral form on the displayed generators; this is
the content of \cite[Section~3]{GN} and the manuscript's
exterior-square model of the automorphic group.
\end{proof}

\begin{remark}[Embedding $\La_{II}^{2,1}\subset\La^{3,2}$]
\label{rem:II21-in-32}
The primitive sublattice
\[
\La^{2,1}=\HypPlan\oplus[2]
=\Z f_2\oplus\Z f_3\oplus\Z f_{-2}\subset\La^{3,2}
\]
of the body's Weyl-chamber section is the BKM root lattice: discarding the
$(f_1,f_{-1})$-block of $\La^{3,2}\simeq\HypPlan\oplus\HypPlan\oplus
[2]$ leaves $\HypPlan\oplus[2]\simeq\La^{2,1}$.  The type-II
sublattice $\La_{II}^{2,1}=\{mf_2+lf_3+nf_{-2}\mid
m,n\equiv 0\!\pmod 2\}$ is the lattice of even integral type-II
classes, on which the simple-root Cartan
\[
\bigl((\delta_i,\delta_j)\bigr)
=\begin{pmatrix}2&-2&-2\\-2&2&-2\\-2&-2&2\end{pmatrix}
\]
of Proposition~\ref{prop:type-II-cartan} is realized via
$\delta_1=2f_2-f_3$, $\delta_2=2f_{-2}-f_3$, $\delta_3=f_3$.  Every
isometry of $\La^{2,1}$ extends to an isometry of $\La^{3,2}$ by
the identity on $(\La^{2,1})^\perp\subset\La^{3,2}$.
\end{remark}

\subsection{Cross-references and consistency check}
\label{app:lattice-cross-refs}

The lattices recorded above sit in the manuscript as follows.

\begin{itemize}
\item $U=\mathrm{II}_{1,1}$: every $U$-summand of
$\La^{3,2}$, $\La_{K3}$, and $\widetilde H(S,\Z)$.  The relation
$\HypPlan=U(-1)$ packages the body's negated convention.
\item $\La_{II}^{2,1}\simeq U(-1)\oplus\langle 2\rangle\simeq
U\oplus\langle-2\rangle$ (sign flip on the rank-$2$ hyperbolic block):
the BKM root lattice on which the type-II simple-root Cartan
$2(\Id_3-\J_3)$ and the Weyl vector
$\rho=\frac12(\delta_1+\delta_2+\delta_3)$ are realized.  The
denominator algebra
$\autcor=\mathfrak g_{\Delta_5}$ is graded by $\La_{II}^{2,1}$
(\cite[Theorem~2.1]{GNII}).
\item $\La_{K3}=\mathrm{II}_{3,19}$: K3 cohomology with
intersection pairing.  Polarized K3 sublattices $\La_S$ embed
primitively into $\La_{K3}$.
\item $\La_S$: signature $(1,\rho-1)$ for $\rho\le 19$ (the
algebraic Picard rank).  Used in
Definition~\ref{def:mukai-lattice} as the K3 factor of
$\Gamma_X^{\mathrm{form}}$ and in the lattice-polarization setup of
the body.
\item $\widetilde H(S,\Z)=\mathrm{II}_{4,20}$: the Mukai lattice on
which the D6--D2--D0 Mukai vectors $v_X(\mathcal I_Y)$ live; the
Gram-index map $\Pi_X$ is computed via this pairing
(Theorem~\ref{thm:d6-d2-d0-dictionary-appendixC}).
\item $\Gamma_X^{\mathrm{form}}=\La_S\oplus\La_S$: the formal
additive dyonic charge lattice of $X=S\times E$, on which the
quadratic Gram-shadow projection $\Pi_X$ and the symmetric bilinear
$2$-cocycle $B$ are defined
(Definition~\ref{def:mukai-lattice}).
\item $\La^{3,2}\simeq U(-1)\oplus U(-1)\oplus\langle 2\rangle$:
the signature-$(3,2)$ orthogonal lattice that bridges
$\PSp_4(\Z)$ to $\SO_+(3,2)$ via the exterior-square
representation; on this lattice the Borcherds product expansion of
$\Delta_5$ is written.
\end{itemize}

\begin{verification}[End-to-end computation check]
\label{ver:lattice-script-check}
\textnormal{[computed, $\gamma$]}
The script \texttt{compute/verify\_lattice.py} verifies, in one
self-contained $\Q$-arithmetic computation:
(i) the rank-$5$ Gram matrix $G_{3,2}$ in the
$(f_1,f_2,f_3,f_{-2},f_{-1})$-basis (lines~82--89);
(ii) the orthogonality of the symplectic invariant
$q_0=e_1\wedge e_3+e_2\wedge e_4$ to the basis (line~90);
(iii) the type-II simple-root Cartan
$2(\Id_3-\J_3)$ at lines~92--97;
(iv) the Weyl vector $\rho=(1,-1/2,1)$ in $(f_2,f_3,f_{-2})$
coordinates and the $(\rho,\delta_i)=-1$ relation
(lines~99--101);
(v) the type-II reflection action and Cartan invariance under each
$s_{\delta_i}$ (lines~103--108).
The \texttt{All lattice checks passed.} line printed at the end of the
script's output confirms agreement of the Gram, Cartan, and
Weyl-vector data of \S\ref{app:lattice-II21} and \S\ref{app:lattice-32}
with the $\Q$-arithmetic computation.
\end{verification}

\begin{remark}[Discriminant convention; sign of $\La^{3,2}$]
\label{rem:discriminant-conventions}
In the body convention $\HypPlan\oplus\HypPlan\oplus[2]$, both
hyperbolic-plane summands have Gram determinant $-1$ and the $[2]$
summand contributes $+2$, so
$\operatorname{disc}\La^{3,2}=(-1)^2\cdot 2=+2$.  In the standard
convention $U\oplus U\oplus\langle-2\rangle$, the two $U$-summands
contribute determinant $-1$ each and $\langle-2\rangle$ contributes
$-2$, so $\operatorname{disc}\La^{3,2}=(-1)^2\cdot(-2)=-2$.  The two
presentations differ by an overall sign on the rank-$2$ hyperbolic
blocks; both carry an order-$2$ discriminant group with single
nontrivial generator $f_3/2\in(\La^{3,2})^*/\La^{3,2}$.  The same
analysis applies to $\La_{II}^{2,1}$: discriminant $-2$ (body) or
$+2$ (standard), discriminant group $\Z/2$ generated by $h/2$.  In
every case the order of the discriminant group is $2$ and the
type-II divisibility statement
$(\delta,\La^{2,1})\in\{\Z,2\Z\}$ depends on the order, not the sign
of the determinant.
\end{remark}

\begin{remark}[Cited primary literature]
\label{rem:lattice-primary-lit}
The primary references behind every load-bearing lattice claim above
are: Conway--Sloane 1999~\S2.5,~\S4.8,~\S15.5,~\S16.4 (classification
of even unimodular indefinite lattices, $E_8$ root system,
discriminant-form gluing, K3 lattice as $\mathrm{II}_{3,19}$);
Nikulin 1980~Theorem~1.1.2,~Theorem~1.4.1 (= \cite{NIKULIN}
Theorem~1.4.1 in the K3-paper version),
Theorem~1.6.1,~Theorem~1.14.4 (signature classification, reflection
group of square-$2$ vectors, primitive embeddings, discriminant-form
uniqueness); Mukai 1984~Theorem~1.5 (Mukai pairing, signature
$(4,20)$), Mukai 1987~Theorem~3 ($\sqrt{\td(S)}=(1,0,1)$ on K3);
Gritsenko--Nikulin 1998 Part~II~Theorem~2.1, Section~5.1.1, case
$(t=1,II,\overline 0)$ (BKM root datum for $\autcor$, Weyl vector,
type-II simple-root Cartan).
\end{remark}
